% =====================================================================
% 06_partie_d.tex — sous-programmes, blocs anonymes et méthodes
% =====================================================================
\chapter{Partie D : sous-programmes et méthodes}
\label{chap:partieD1}

\section{Objectifs de la programmation PL/SQL}

La partie D transforme le schéma en un système actif. Elle ajoute des
sous-programmes catalogués, des blocs anonymes, des méthodes de types objet et
des triggers. Ce chapitre couvre les deux premières catégories, ainsi que les
méthodes. Les triggers, les transactions et la clôture de compte sont traités
au chapitre~\ref{chap:partieD2}, car ils modifient l'état de la base et
nécessitent une chronologie précise.

\begin{objectifbox}[Trois niveaux de programmation]
Le projet distingue trois niveaux : les sous-programmes catalogués, appelables
directement en SQL*Plus ; les blocs anonymes, utilisés pour les démonstrations
ponctuelles ; et les méthodes de types objet, qui encapsulent les règles dans
les objets eux-mêmes. Cette séparation évite de mélanger la logique métier et
les scripts de test.
\end{objectifbox}

\section{Sous-programmes catalogués}

Neuf sous-programmes catalogués sont utilisés dans le projet. Trois sont des
fonctions de lecture, quatre sont des procédures de gestion et deux proviennent
de la partie B. Le tableau~\ref{tab:sous-programmes} rappelle leur signature,
leur rôle et leur comportement transactionnel.

\begin{table}[htbp]
  \centering
  \small
  \caption{Sous-programmes catalogués.}
  \label{tab:sous-programmes}
  \footnotesize
  \begin{tabularx}{\textwidth}{@{}L{4.0cm}L{5.6cm}X@{}}
    \toprule
    \entete{Sous-programme} & \entete{Signature} & \entete{Rôle et transaction} \\
    \midrule
    \code{nb\_comptes} &
    \code{FUNCTION nb\_comptes RETURN NUMBER} &
    Nombre total de comptes (CC + CE + CT). Lecture seule. \\
    \code{solde\_total} &
    \code{FUNCTION solde\_total(paramcli IN NUMBER)} &
    Somme des soldes d'un client. Lecture seule. \\
    \code{ville\_client} &
    \code{FUNCTION ville\_client(paramcli IN NUMBER)} &
    Ville d'un client. Lecture seule. \\
    \code{ajout\_client} &
    \code{PROCEDURE ajout\_client(...)} &
    Ajout complet d'un client, parrain optionnel. Pas de \code{COMMIT}. \\
    \code{virement} &
    \code{PROCEDURE virement(cptsrc, cptdest, montant, cli)} &
    Virement atomique entre deux comptes courants, commission
    inter-agences. Pas de \code{COMMIT}. \\
    \code{maj\_statuts\_dormants} &
    \code{PROCEDURE maj\_statuts\_dormants} &
    Passe en \code{DORMANT} les comptes courants sans opération depuis
    douze mois. Pas de \code{COMMIT}. \\
    \code{scan\_echeances} &
    \code{PROCEDURE scan\_echeances} &
    Crée les alertes d'échéance pour les comptes à terme. \code{COMMIT}. \\
    \code{change\_Portable} &
    \code{PROCEDURE change\_Portable(...)} &
    Modifie le deuxième téléphone. \code{COMMIT} en succès,
    \code{ROLLBACK} en erreur. \\
    \code{drop\_client} &
    \code{PROCEDURE drop\_client(...)} &
    Supprime un client et ses dépendances. \code{COMMIT} ou
    \code{ROLLBACK}. \\
    \bottomrule
  \end{tabularx}
\end{table}

\begin{clebox}[Pourquoi \code{virement} ne fait-il pas \code{COMMIT} ?]
Un virement peut faire partie d'une transaction plus large. Le \code{COMMIT}
appartient à l'appelant. La procédure utilise un \code{SAVEPOINT} et un
\code{ROLLBACK TO} pour garantir l'atomicité des deux écritures, sans imposer
la validation. Cette conception est essentielle pour tester le rollback
global.
\end{clebox}

\section{Le virement atomique}

La procédure \code{virement} vérifie que le montant est strictement positif et
que les comptes source et destination sont différents. Elle récupère les
\code{REF} des comptes et du client, calcule la commission inter-agences, puis
insère un débit et un crédit. En cas d'erreur, elle revient au point de
sauvegarde : les deux mouvements sont annulés ensemble.

\begin{lstlisting}[caption={Procédure \code{virement} : atomicité par \code{SAVEPOINT}.},label={lst:virement}]
CREATE OR REPLACE PROCEDURE virement (
  cptsrc  IN VARCHAR2,
  cptdest IN VARCHAR2,
  montant IN NUMBER,
  cli     IN NUMBER
) IS
  v_ref_src  REF cptCourant_type;
  v_ref_dest REF cptCourant_type;
  v_ref_cli  REF client_type;
  v_ag_src   VARCHAR2(6);
  v_ag_dest  VARCHAR2(6);
  v_comm     NUMBER(12,2) := 0;
BEGIN
  SAVEPOINT av_virement;

  IF montant <= 0 THEN
    RAISE_APPLICATION_ERROR(-20020,
      'Le montant doit etre strictement positif.');
  END IF;
  IF cptsrc = cptdest THEN
    RAISE_APPLICATION_ERROR(-20021,
      'Comptes source et destination identiques.');
  END IF;

  SELECT REF(x) INTO v_ref_src
    FROM cpt_courant x WHERE x.num_compte = cptsrc;
  SELECT REF(x) INTO v_ref_dest
    FROM cpt_courant x WHERE x.num_compte = cptdest;
  SELECT REF(c) INTO v_ref_cli
    FROM client c WHERE c.num_client = cli;

  SELECT x.ref_agence.code INTO v_ag_src
    FROM cpt_courant x WHERE x.num_compte = cptsrc;
  SELECT x.ref_agence.code INTO v_ag_dest
    FROM cpt_courant x WHERE x.num_compte = cptdest;
  IF v_ag_src <> v_ag_dest THEN
    v_comm := GREATEST(ROUND(montant * 0.01, 2), 100);
  END IF;

  INSERT INTO operation (id_op, ref_client, ref_compte, date_op,
                         montant, type_op, commission)
  VALUES (seq_operation.NEXTVAL, v_ref_cli, v_ref_src, TRUNC(SYSDATE),
          -montant, 'VIREMENT', v_comm);

  INSERT INTO operation (id_op, ref_client, ref_compte, date_op,
                         montant, type_op, commission)
  VALUES (seq_operation.NEXTVAL, v_ref_cli, v_ref_dest, TRUNC(SYSDATE),
          montant, 'VIREMENT', 0);

  DBMS_OUTPUT.PUT_LINE('virement OK : '||montant||' DA de '||cptsrc
                       ||' vers '||cptdest||' (commission='||v_comm||')');
EXCEPTION
  WHEN OTHERS THEN
    ROLLBACK TO av_virement;
    DBMS_OUTPUT.PUT_LINE('virement ANNULE : '||SQLERRM);
END;
/
\end{lstlisting}

\begin{remarquebox}[Commission inter-agences]
La commission vaut 1~\% du montant, avec un minimum de 100~DA, uniquement si
les agences source et destination diffèrent. Elle est portée par l'opération de
débit et reste nulle sur le crédit. Ce choix évite de créer une troisième
opération de commission et garde la trace du prélèvement sur le compte source.
\end{remarquebox}

\section{Blocs anonymes D7}

Les blocs anonymes illustrent des traitements ponctuels : lecture d'un élément
de \code{VARRAY}, parcours d'un curseur, curseur paramétré, extension d'une
collection et calcul d'intérêts. Ils ne créent pas d'objet persistant, mais ils
montrent la maîtrise des structures PL/SQL.

\begin{table}[htbp]
  \centering
  \small
  \caption{Blocs anonymes D7a à D7e.}
  \label{tab:blocs-anonymes}
  \begin{tabularx}{\textwidth}{@{}C{0.9cm}L{4.7cm}L{4.4cm}X@{}}
    \toprule
    \entete{Bloc} & \entete{Consigne} & \entete{Technique} &
    \entete{Résultat} \\
    \midrule
    7a & Téléphone travail (2\up{e} élément) du client 3 &
    \code{v\_tel(2)} + test \code{NULL} & \code{0777-11-22-33}. \\
    7b & Même information pour tous les titulaires de CC &
    Curseur + \code{FOR} &
    Affiche les 6 clients, certains « non renseigné ». \\
    7c & Curseur paramétré par un seuil, titulaire, solde, nombre de
    signataires & Curseur paramétré + \code{NO\_DATA\_FOUND} +
    \code{aucun\_compte} &
    Seuil 100\,000 : CC2, CC3, CC5, CC6 ; seuil $-1\,000\,000$ :
    \code{aucun\_compte}. \\
    7d & Ajouter un 3\up{e} téléphone avec \code{EXTEND}, \code{LIMIT},
    \code{COUNT}, \code{LAST} & Client 98 temporaire &
    \code{COUNT} 2 $\rightarrow$ 3 ; 4\up{e} refusé ; client 98 supprimé
    ensuite. \\
    7e & Total des intérêts annuels des CE d'une agence &
    Curseur + \code{ROUND(solde * taux / 100, 2)} &
    A1 : 10\,150~DA. \\
    \bottomrule
  \end{tabularx}
\end{table}

Le bloc 7e est particulièrement utile pour vérifier la cohérence des calculs
d'intérêts. Pour l'agence A1, il additionne les intérêts de CE1 et CE4 :
$150\,000 \times 2,5\,\% = 3\,750$ et
$320\,000 \times 2\,\% = 6\,400$, soit 10\,150~DA. CE4 existe encore à ce
moment ; il est supprimé plus tard par \code{fermeCompte('CE4')}.

\begin{lstlisting}[caption={Bloc 7e : total des intérêts annuels de l'agence A1.},label={lst:bloc-7e}]
DECLARE
  CURSOR c_epargne (p_agence VARCHAR2) IS
    SELECT ce.num_compte, ce.solde, ce.taux_interet
      FROM cpt_epargne ce, agence a
     WHERE ce.ref_agence = REF(a)
       AND a.code = p_agence
     ORDER BY ce.num_compte;
  v_interet NUMBER;
  v_total   NUMBER := 0;
BEGIN
  DBMS_OUTPUT.PUT_LINE('Interets annuels des comptes d epargne de A1 :');
  FOR r IN c_epargne('A1') LOOP
    v_interet := ROUND(r.solde * r.taux_interet / 100, 2);
    DBMS_OUTPUT.PUT_LINE('  '||r.num_compte
                         ||' : solde='||r.solde
                         ||' taux='||r.taux_interet||'%'
                         ||' interet='||v_interet);
    v_total := v_total + v_interet;
  END LOOP;
  DBMS_OUTPUT.PUT_LINE('  TOTAL ANNUEL = '||v_total);
END;
/
\end{lstlisting}

\section{Méthodes des types objet}

Les méthodes sont ajoutées aux types par \code{ALTER TYPE ... ADD MEMBER},
puis implantées dans les corps de types. Chaque \code{ALTER TYPE} est protégé
contre l'erreur \code{ORA-22324} afin que le script puisse être rejoué sans
recréer la partie A. Oracle n'accepte qu'un seul membre par instruction
\code{ALTER TYPE}.

\begin{table}[htbp]
  \centering
  \small
  \caption{Méthodes des types objet et résultats clés.}
  \label{tab:methodes}
  \begin{tabularx}{\textwidth}{@{}L{3.4cm}L{4.7cm}X@{}}
    \toprule
    \entete{Type} & \entete{Méthode} & \entete{Rôle et résultat clé} \\
    \midrule
    \code{client\_type} & \code{nbCepargne} &
    Nombre de CE du client : 2, 1, 1, 1, 0, 0. \\
    \code{client\_type} & \code{soldeTotal} &
    Somme CC + CE + CT : 3\,348\,500 ; 300\,000 ; 168\,000 ; 830\,000 ;
    0 ; 40\,000. \\
    \code{client\_type} & \code{supprimeTel(indice)} &
    Supprime un téléphone. Client 4 : indice 2 supprimé puis restauré ;
    client 5 : indice 2 non renseigné. \\
    \code{client\_type} & \code{afficheInterets} &
    Affiche les intérêts des CE. Client 1 : 18\,750~DA. \\
    \code{cptCourant\_type} & \code{nbSignataire} /
    \code{nbSignataire(droit)} &
    Nombre de signataires, éventuellement par droit. Surcharge. \\
    \code{cptCourant\_type} & \code{estTitulaire(cli)} &
    Booléen : le client est-il titulaire ? Testé en PL/SQL. \\
    \code{cptCourant\_type} & \code{estSignataire(cli)} &
    Booléen : le client est-il signataire ? Testé en PL/SQL. \\
    \code{cptCourant\_type} & \code{supprimeSign(cli)} &
    Retire un signataire. CC7 : 3 $\rightarrow$ 2, puis restauration. \\
    \code{cptCourant\_type} & \code{peutRetirer(cli, montant)} &
    Vérifie les droits et le découvert.
    \code{peutRetirer(4, -1000) = TRUE}. \\
    \code{cptCourant\_type} & \code{passeMvt(cli, somme)} &
    Insère une opération et met à jour le solde. CC7 $+10\,000 \rightarrow
    195\,000$ ; $-5\,000 \rightarrow 180\,000$ (tests séparés). \\
    \code{cptCourant\_type} & \code{fraisTenue} &
    Frais de tenue : 1\,200~DA pour un CC, 0 pour CE/CT. \\
    \code{cptEpargne\_type} & \code{calculeInterets} &
    Solde + intérêts. CE5 : 515\,000. \\
    \code{cptEpargne\_type} & \code{sortieAnticipee} &
    Solde avec pénalité de 5~\%. CE5 : 475\,000. \\
    \code{cptTerme\_type} & \code{calculeInterets} &
    Solde + intérêts proratisés. CT1 : 206\,400. \\
    \code{cptTerme\_type} & \code{sortieAnticipee} &
    Solde avec pénalité de 5~\%. CT1 : 190\,000. \\
    \code{compte\_type} & \code{fermeCompte(cpt)} (statique) &
    Clôture un compte. CE4 : 320\,000 ; CC1 : 2\,498\,500 et cinq
    opérations. \\
    \bottomrule
  \end{tabularx}
\end{table}

\begin{clebox}[Pourquoi \code{calculeInterets} retourne-t-il le solde plus les intérêts ?]
La méthode simule le montant disponible à l'échéance : pour l'épargne,
\code{solde * (1 + taux/100)} ; pour le terme,
\code{solde * (1 + taux/100 * durée/12)}. Les intérêts seuls sont visibles
dans \code{afficheInterets} : 3\,750 + 15\,000 = 18\,750~DA pour le client 1.
\end{clebox}

\begin{lstlisting}[caption={Calcul des intérêts et de la sortie anticipée.},label={lst:interets}]
MEMBER FUNCTION calculeInterets RETURN NUMBER IS
BEGIN
  RETURN ROUND(SELF.solde * (1 + SELF.taux_interet / 100), 2);
END;

MEMBER FUNCTION sortieAnticipee RETURN NUMBER IS
BEGIN
  RETURN ROUND(SELF.solde * 0.95, 2);
END;
\end{lstlisting}

Pour un compte à terme, la durée intervient dans le calcul :
\code{solde * (1 + (taux\_garanti / 100) * (duree\_mois / 12))}. La pénalité de
sortie anticipée reste de 5~\%, quelle que soit la durée restante. Ce choix
simplifie la règle métier et la rend facile à expliquer.

\section{Clôture de compte}

La procédure statique \code{fermeCompte} cherche le compte dans
\code{cpt\_courant}, puis dans \code{cpt\_epargne}, puis dans
\code{cpt\_terme}. Elle affiche le solde avant clôture, supprime les opérations
si c'est un compte courant, puis supprime le compte. Elle valide par
\code{COMMIT}. Si le compte n'existe pas, elle affiche un message et annule.

\begin{attentionbox}[Procédure statique]
\code{fermeCompte} est déclarée \code{STATIC} dans \code{compte\_type} : elle
est appelée sur le type, et non sur une instance. L'appel
\code{compte\_type.fermeCompte('CC1')} ne dépend donc pas d'un objet
\code{compte\_type} préalable. C'est cohérent avec une opération de clôture
identifiée par un numéro de compte.
\end{attentionbox}

\section{Tests et résultats}

Les tests des sous-programmes, des blocs anonymes et des méthodes sont
exécutés avant les triggers, sauf pour \code{passeMvt} et \code{fermeCompte},
qui nécessitent les triggers et sont testés plus tard. Le
tableau~\ref{tab:resultats-D1} récapitule les résultats attendus.

\begin{table}[htbp]
  \centering
  \small
  \caption{Résultats des tests D1, D7 et D8.}
  \label{tab:resultats-D1}
  \begin{tabularx}{\textwidth}{@{}L{5.2cm}L{4.6cm}X@{}}
    \toprule
    \entete{Test} & \entete{Commande ou méthode} & \entete{Résultat attendu} \\
    \midrule
    Nombre de comptes & \code{SELECT nb\_comptes() FROM dual;} &
    Somme des CC, CE et CT. \\
    Solde total client 1 & \code{SELECT solde\_total(1) FROM dual;} &
    3\,348\,500. \\
    Ville client 1 & \code{SELECT ville\_client(1) FROM dual;} &
    Béjaïa. \\
    Ajout client 97 & \code{ajout\_client(97, ...)} &
    Client ajouté, puis supprimé. \\
    Nb épargne par client & \code{c.nbCepargne()} &
    2, 1, 1, 1, 0, 0. \\
    Solde total par client & \code{c.soldeTotal()} &
    3\,348\,500 ; 300\,000 ; 168\,000 ; 830\,000 ; 0 ; 40\,000. \\
    Intérêts client 1 & \code{v\_client.\allowbreak afficheInterets()} &
    18\,750~DA. \\
    Nombre de signataires & \code{cc.nbSignataire()} &
    CC1:1, CC2:1, CC3:2, CC4:3, CC5:0, CC6:0, CC7:3. \\
    Test de retrait & \code{peutRetirer(4, -1000)} & \code{TRUE}. \\
    Épargne CE5 & \code{calculeInterets()} / \code{sortieAnticipee()} &
    515\,000 / 475\,000. \\
    Terme CT1 & \code{calculeInterets()} / \code{sortieAnticipee()} &
    206\,400 / 190\,000. \\
    \bottomrule
  \end{tabularx}
\end{table}

\section{Pièges de la programmation PL/SQL}

\begin{table}[htbp]
  \centering
  \small
  \caption{Pièges de la partie D1.}
  \label{tab:pieges-D1}
  \begin{tabularx}{\textwidth}{@{}L{5.0cm}L{4.8cm}X@{}}
    \toprule
    \entete{Piège} & \entete{Symptôme} & \entete{Solution} \\
    \midrule
    Ajouter un membre déjà présent & \code{ORA-22324}. &
    Protéger l'\code{ALTER TYPE} par une gestion d'exception ou tester
    \code{user\_type\_attrs}. \\
    \code{ALTER TYPE} multiple & Erreur de syntaxe. &
    Un seul \code{ADD MEMBER} par instruction. \\
    \code{COMMIT} dans une procédure réutilisable &
    Le rollback global ne fonctionne plus. &
    Laisser le \code{COMMIT} à l'appelant, sauf opération autonome. \\
    \code{BOOLEAN} en SQL & \code{ORA-00902}. &
    Utiliser \code{BOOLEAN} uniquement en PL/SQL, ou convertir en
    \code{CHAR(1)}. \\
    \code{VALUE} sur une table objet & Objet non récupéré. &
    Utiliser \code{SELECT VALUE(c) INTO ...}. \\
    \bottomrule
  \end{tabularx}
\end{table}

\section{Conclusion de la partie D1}

Les sous-programmes et les méthodes donnent au schéma une réelle capacité de
traitement. Ils encapsulent les calculs de soldes, d'intérêts, de droits et de
clôture. Le choix de ne pas valider dans \code{virement} ou
\code{ajout\_client} prépare les tests transactionnels du chapitre suivant. Les
méthodes polymorphes, en particulier \code{fraisTenue} et
\code{calculeInterets}, montrent que le comportement dépend du type réel de
l'objet et non d'un test conditionnel externe.
